% Generated draft: requires mathematical and source review.
\documentclass[11pt]{article}
\usepackage[margin=25mm]{geometry}
\usepackage{fontspec}
\setmainfont{Latin Modern Roman}
\usepackage{amsmath,amssymb,mathtools,unicode-math}
\setmathfont{Latin Modern Math}
\usepackage{enumitem,needspace,xurl,hyperref}
\setlength{\parindent}{0pt}
\setlength{\parskip}{5pt}
\setlength{\emergencystretch}{3em}
\newcommand{\N}{\mathbb{N}}
\newcommand{\Z}{\mathbb{Z}}
\newcommand{\Q}{\mathbb{Q}}
\newcommand{\R}{\mathbb{R}}
\newcommand{\C}{\mathbb{C}}
\newcommand{\F}{\mathbb{F}}
\newcommand{\A}{\mathbb{A}}
\newcommand{\PP}{\mathbb{P}}
\newcommand{\E}{\mathbb{E}}
\newcommand{\eps}{\varepsilon}
\newcommand{\dd}{\,\mathrm{d}}
\begin{document}
\section*{UnsolvedMath ID: 6000016}
\textbf{Source problem code:} catalogue entry
\subsection{Definitions and mathematical statement}
A Hessian manifold in the usual positive-definite convention is a smooth manifold $M$ with a flat torsion-free affine connection $\nabla$ and a Riemannian metric $g$ locally of the form
\[
 g_{ij}=\frac{\partial^2\varphi}{\partial x^i\partial x^j}
\]
in $\nabla$-affine coordinates, for local strictly convex potentials $\varphi$. Equivalently $\nabla g$ is totally symmetric. These potentials need only be local; no globally defined potential is required.

Let $(M,\nabla_0,g_0)$ be compact. For a chosen deformation topology on affine structures on the same underlying smooth manifold, the source asks whether there is a neighborhood $U$ of $\nabla_0$ such that
\[
 \forall\nabla\in U\text{ flat and torsion free},\qquad
 \exists g\text{ positive definite with }\nabla g
             \text{ totally symmetric}.
\]
Equivalently, the question asks whether admitting some Hessian metric is locally stable under the intended small affine deformations at every compact Hessian structure.

The source does not specify the topology, whether affine structures are marked or taken modulo diffeomorphisms, or whether deformations are arbitrary nearby structures or prescribed parameterized families. Its existence wording requires a Hessian metric for each deformed structure; it does not explicitly require a choice varying smoothly, a metric close to $g_0$, preservation of a cohomology class, or any deformation of $g_0$ itself. These stronger conditions are not added.
\subsection{Short English statement}
Does every sufficiently small change of the flat affine structure of a compact Hessian manifold still support some Hessian metric? The metric may be newly chosen; the source does not demand a smooth family of metrics or a fixed potential class.

The supplied summary reports nontrivial affine deformations on tori and describes Hessian-metric existence as constrained by cohomological or convex-cone conditions. It calls the stability question open or only partially addressed and says no definitive answer was verified. Its phrase ``not a generic open property in general'' is not backed by a precise obstruction or counterexample in the record and is not treated as a resolution. Provenance marks current status of the 1998 problem for review.
\subsection{Sources}
\begin{itemize}
\item Local UnsolvedMath registry, permanent ID 6000016, source code catalogue entry, ``Stability of Hessian Metrics,'' proposed by Satoru Shimizu, 1998. Dataset curation and metadata: Ulam AI and UnsolvedMath Contributors, CC BY 4.0.
\item Furuhata, Matsuzoe, and Urakawa, \emph{Open Problems in Affine Differential Geometry and Related Topics} (1998), item 5(b), scanned PDF pages 1--2, \url{https://www.jstage.jst.go.jp/article/iis/4/2/4_2_125/_pdf/-char/en}; manual OCR from page images, accessed 2026-07-29.
\item The research summary's torus-deformation and cohomological/convex-cone descriptions have no precise theorem statements or separate primary references.
\item Source-specific provenance: the article presents this as an open problem and current status NEEDS\_REVIEW. The J-STAGE article page marks the article CC BY 4.0, retained in addition to dataset attribution.
\end{itemize}
\end{document}
